% GENERATED FILE. Edit canonical lessons, then run npm run generate:books.
% canonical-source-hash: fnv1a64:b63ef44ce1b6e669
% canonical-lessons: TA-C362-araycci, TA-C362-ayvu, TA-C362-taravu, TA-C362-pakuppayvu, TA-C362-alacu, TA-R362-review-355-362

\chapter{Doing Research}
\label{ch:ta-doing-research}
\hlchaptermodality{362}

\begin{chapteropening}
\textbf{By the end of this chapter:} \emph{I can talk about research, data and analysis.}

Complete the last lesson of chapter 362: I can talk about research, data and analysis.
\end{chapteropening}

\section[\texorpdfstring{\ta{ஆராய்ச்சி} (ārāycci)}{ஆராய்ச்சி (ārāycci)}]{\ta{ஆராய்ச்சி} (ārāycci) — research}

\label{lesson:TA-C362-araycci}



\begin{quote}
Before the new one: say the Tamil for a good match, then the Tamil for a connection.
\end{quote}

\subsection*{You'll want to know: \ta{ஆராய்ச்சி}}
\textbf{\ta{ஆராய்ச்சி}} — \emph{ārāycci} — \textquotedblleft{}research\textquotedblright{}.

\begin{quote}
\ta{என்} \ta{அக்கா} \ta{செடிகளைப்} \ta{பற்றி} \ta{ஆராய்ச்சி} \ta{செய்கிறாள்}.
\end{quote}

\begin{quote}
\emph{eṉ akkā ceṭikaḷaip paṟṟi ārāycci ceykiṟāḷ.}
\end{quote}

\begin{quote}
My elder sister does research on plants.
\end{quote}

\ta{ஆராய்ச்சி} \ta{செய்}, to do research. Its verb is \ta{ஆராய்}, to investigate.

\begin{cousinweb}
\textbf{\ta{ஆராய்ச்சி}} (\emph{ārāycci}) is native Dravidian: DEDR 377 lists Tamil \emph{ārāy}, \textquotedblleft{}to investigate, examine, seek\textquotedblright{}, and \emph{ārāycci}, \textquotedblleft{}research, investigation\textquotedblright{}.
\end{cousinweb}

\begin{grammarlens}[title={What you've built}]
One more word for giving reasons and opinions.
\end{grammarlens}

\subsection*{Your turn}
\begin{itemize}
\raggedright
  \item \emph{Say it:} \emph{ārāycci}
  \item \emph{Say it:} \emph{ārāycci}, once more
  \item \emph{Say it:} \emph{eṉ akkā ceṭikaḷaip paṟṟi ārāycci ceykiṟāḷ.} — My elder sister does research on plants.
  \item \emph{Recall:} say the Tamil for a good match, then the Tamil for a connection, then say \emph{ārāycci} again
\end{itemize}

\begin{tcolorbox}[breakable,colback=teal!4,colframe=teal!35!black,title={Before you move on}]
What does \ta{ஆராய்ச்சி} mean? (\textquotedblleft{}Research\textquotedblright{}.) Say it once more.
\end{tcolorbox}
\section[\texorpdfstring{\ta{ஆய்வு} (āyvu)}{ஆய்வு (āyvu)}]{\ta{ஆய்வு} (āyvu) — a study; an inspection}

\label{lesson:TA-C362-ayvu}



\begin{quote}
Before the new one: say the Tamil for a connection, then the Tamil for research.
\end{quote}

\subsection*{You'll want to know: \ta{ஆய்வு}}
\textbf{\ta{ஆய்வு}} — \emph{āyvu} — \textquotedblleft{}a study; an inspection\textquotedblright{}.

\begin{quote}
\ta{இந்த} \ta{ஆய்வு} \ta{மூன்று} \ta{ஆண்டுகள்} \ta{நடந்தது}.
\end{quote}

\begin{quote}
\emph{inta āyvu mūṉṟu āṇṭukaḷ naṭantatu.}
\end{quote}

\begin{quote}
This study went on for three years.
\end{quote}

\ta{ஆய்வு} \ta{செய்}, to study, to inspect. \ta{ஆய்வாளர்}, a researcher.

\begin{cousinweb}
\textbf{\ta{ஆய்வு}} (\emph{āyvu}) is native Dravidian: DEDR 363 lists Tamil \emph{āy}, \textquotedblleft{}to search, examine, select\textquotedblright{}, and \emph{āyvu}, \textquotedblleft{}investigation, research\textquotedblright{}.
\end{cousinweb}

\begin{grammarlens}[title={What you've built}]
One more word for giving reasons and opinions.
\end{grammarlens}

\subsection*{Your turn}
\begin{itemize}
\raggedright
  \item \emph{Say it:} \emph{āyvu}
  \item \emph{Say it:} \emph{āyvu}, once more
  \item \emph{Say it:} \emph{inta āyvu mūṉṟu āṇṭukaḷ naṭantatu.} — This study went on for three years.
  \item \emph{Recall:} say the Tamil for a connection, then the Tamil for research, then say \emph{āyvu} again
\end{itemize}

\begin{tcolorbox}[breakable,colback=teal!4,colframe=teal!35!black,title={Before you move on}]
What does \ta{ஆய்வு} mean? (\textquotedblleft{}A study; an inspection\textquotedblright{}.) Say it once more.
\end{tcolorbox}
\section[\texorpdfstring{\ta{தரவு} (taravu)}{தரவு (taravu)}]{\ta{தரவு} (taravu) — data}

\label{lesson:TA-C362-taravu}



\begin{quote}
Before the new one: say the Tamil for research, then the Tamil for a study.
\end{quote}

\subsection*{You'll want to know: \ta{தரவு}}
\textbf{\ta{தரவு}} — \emph{taravu} — \textquotedblleft{}data\textquotedblright{}.

\begin{quote}
\ta{இந்தத்} \ta{தரவுகளை} \ta{எங்கிருந்து} \ta{பெற்றீர்கள்}?
\end{quote}

\begin{quote}
\emph{intat taravukaḷai eṅkiruntu peṟṟīrkaḷ?}
\end{quote}

\begin{quote}
Where did you get this data?
\end{quote}

Plural: \ta{தரவுகள்}, \emph{taravukaḷ}.

Plural \ta{தரவுகள்}, the pieces of data.

\begin{cousinweb}
\textbf{\ta{தரவு}} (\emph{taravu}) is a modern Tamil coinage for \textquotedblleft{}data\textquotedblright{}, usually read as built on native \emph{taru}, \textquotedblleft{}to give\textquotedblright{} (DEDR 3098): \textquotedblleft{}what is given\textquotedblright{}, which is also what Latin \emph{data} means.
\end{cousinweb}

\begin{grammarlens}[title={What you've built}]
One more word for giving reasons and opinions.
\end{grammarlens}

\subsection*{Your turn}
\begin{itemize}
\raggedright
  \item \emph{Say it:} \emph{taravu}
  \item \emph{Say it:} \emph{taravu}, once more
  \item \emph{Say it:} \emph{intat taravukaḷai eṅkiruntu peṟṟīrkaḷ?} — Where did you get this data?
  \item \emph{Recall:} say the Tamil for research, then the Tamil for a study, then say \emph{taravu} again
\end{itemize}

\begin{tcolorbox}[breakable,colback=teal!4,colframe=teal!35!black,title={Before you move on}]
What does \ta{தரவு} mean? (\textquotedblleft{}Data\textquotedblright{}.) Say it once more.
\end{tcolorbox}
\section[\texorpdfstring{\ta{பகுப்பாய்வு} (pakuppāyvu)}{பகுப்பாய்வு (pakuppāyvu)}]{\ta{பகுப்பாய்வு} (pakuppāyvu) — an analysis}

\label{lesson:TA-C362-pakuppayvu}



\begin{quote}
Before the new one: say the Tamil for a study, then the Tamil for data.
\end{quote}

\subsection*{You'll want to know: \ta{பகுப்பாய்வு}}
\textbf{\ta{பகுப்பாய்வு}} — \emph{pakuppāyvu} — \textquotedblleft{}an analysis\textquotedblright{}.

\begin{quote}
\ta{ஆசிரியர்} \ta{தேர்வு} \ta{முடிவுகளைப்} \ta{பகுப்பாய்வு} \ta{செய்தார்}.
\end{quote}

\begin{quote}
\emph{āciriyar tērvu muṭivukaḷaip pakuppāyvu ceytār.}
\end{quote}

\begin{quote}
The teacher analysed the exam results.
\end{quote}

\ta{பகுப்பாய்வு} \ta{செய்}, to analyse.

\begin{cousinweb}
\textbf{\ta{பகுப்பாய்வு}} (\emph{pakuppāyvu}) is a modern native coinage: \emph{pakuppu}, \textquotedblleft{}division\textquotedblright{}, from \emph{paku}, \textquotedblleft{}to divide\textquotedblright{} (DEDR 3808), and \emph{āyvu}, \textquotedblleft{}examination\textquotedblright{} (DEDR 363): taking apart and examining.
\end{cousinweb}

\begin{grammarlens}[title={What you've built}]
One more word for giving reasons and opinions.
\end{grammarlens}

\subsection*{Your turn}
\begin{itemize}
\raggedright
  \item \emph{Say it:} \emph{pakuppāyvu}
  \item \emph{Say it:} \emph{pakuppāyvu}, once more
  \item \emph{Say it:} \emph{āciriyar tērvu muṭivukaḷaip pakuppāyvu ceytār.} — The teacher analysed the exam results.
  \item \emph{Recall:} say the Tamil for a study, then the Tamil for data, then say \emph{pakuppāyvu} again
\end{itemize}

\begin{tcolorbox}[breakable,colback=teal!4,colframe=teal!35!black,title={Before you move on}]
What does \ta{பகுப்பாய்வு} mean? (\textquotedblleft{}An analysis\textquotedblright{}.) Say it once more.
\end{tcolorbox}
\section[\texorpdfstring{\ta{அலசு} (alacu)}{அலசு (alacu)}]{\ta{அலசு} (alacu) — to analyse closely; to rinse}

\label{lesson:TA-C362-alacu}



\begin{quote}
Before the new one: say the Tamil for data, then the Tamil for an analysis.
\end{quote}

\subsection*{You'll want to know: \ta{அலசு}}
\textbf{\ta{அலசு}} — \emph{alacu} — \textquotedblleft{}to analyse closely; to rinse\textquotedblright{}.

\begin{quote}
\ta{மாணவர்கள்} \ta{இந்தப்} \ta{பிரச்சினையை} \ta{விரிவாக} \ta{அலசினார்கள்}.
\end{quote}

\begin{quote}
\emph{māṇavarkaḷ intap piracciṉaiyai virivāka alaciṉārkaḷ.}
\end{quote}

\begin{quote}
The students went over this problem in detail.
\end{quote}

Past (I …): \ta{அலசினேன்}, \emph{alaciṉēṉ}.

The first sense is to rinse clothes; a question is gone over again and again in the same way.

\begin{cousinweb}
\textbf{\ta{அலசு}} (\emph{alacu}) is native Dravidian: DEDR 246 lists Tamil \emph{alampu}, \emph{alacu}, \textquotedblleft{}to wash, rinse\textquotedblright{}. Going over a question closely is its newer sense.
\end{cousinweb}

\begin{grammarlens}[title={What you've built}]
One more word for giving reasons and opinions.
\end{grammarlens}

\subsection*{Your turn}
\begin{itemize}
\raggedright
  \item \emph{Say it:} \emph{alacu}
  \item \emph{Say it:} \emph{alacu}, once more
  \item \emph{Say it:} \emph{māṇavarkaḷ intap piracciṉaiyai virivāka alaciṉārkaḷ.} — The students went over this problem in detail.
  \item \emph{Recall:} say the Tamil for data, then the Tamil for an analysis, then say \emph{alacu} again
\end{itemize}

\begin{tcolorbox}[breakable,colback=teal!4,colframe=teal!35!black,title={Before you move on}]
What does \ta{அலசு} mean? (\textquotedblleft{}To analyse closely; to rinse\textquotedblright{}.) Say it once more.
\end{tcolorbox}
\section[(dialogue)]{Looking back: from Advice and Proposals to Doing Research}

\label{lesson:TA-R362-review-355-362}



\begin{quote}
Say the Tamil for to recommend and advice.
\end{quote}

\subsection*{Your turn}
Say these aloud:

\begin{itemize}
\raggedright
  \item to recommend, advice, to advise, to consult, to propose
  \item to talk over, negotiations, to interrupt, a compromise, a heated argument
  \item nevertheless, on the other hand, except, likewise, to differ
  \item a theory, a philosophy, an ideology, a policy, to establish
  \item to twist, to exaggerate, to contradict, uncertainty, to reject
  \item fairness, justice, ethics, blame, an excuse
  \item to determine, to look into, a yardstick, a good match, a connection
  \item research, a study, data, an analysis, to analyse closely
\end{itemize}

\begin{tcolorbox}[breakable,colback=teal!4,colframe=teal!35!black,title={Before you move on}]
To recommend? (\textbf{\ta{பரிந்துரை}}, \emph{parinturai}.) Advice? (\textbf{\ta{அறிவுரை}}, \emph{aṟivurai}.) To advise? (\textbf{\ta{அறிவுறுத்து}}, \emph{aṟivuṟuttu}.) To consult? (\textbf{\ta{ஆலோசி}}, \emph{ālōci}.) To propose? (\textbf{\ta{முன்மொழி}}, \emph{muṉmoḻi}.) To talk over? (\textbf{\ta{கலந்துரையாடு}}, \emph{kalanturaiyāṭu}.) Negotiations? (\textbf{\ta{பேச்சுவார்த்தை}}, \emph{pēccuvārttai}.) To interrupt? (\textbf{\ta{குறுக்கிடு}}, \emph{kuṟukkiṭu}.) A compromise? (\textbf{\ta{சமரசம்}}, \emph{camaracam}.) A heated argument? (\textbf{\ta{வாக்குவாதம்}}, \emph{vākkuvātam}.) Nevertheless? (\textbf{\ta{எனினும்}}, \emph{eṉiṉum}.) On the other hand? (\textbf{\ta{மறுபுறம்}}, \emph{maṟupuṟam}.) Except? (\textbf{\ta{தவிர}}, \emph{tavira}.) Likewise? (\textbf{\ta{அதேபோல்}}, \emph{atēpōl}.) To differ? (\textbf{\ta{வேறுபடு}}, \emph{vēṟupaṭu}.) A theory? (\textbf{\ta{கோட்பாடு}}, \emph{kōṭpāṭu}.) A philosophy? (\textbf{\ta{தத்துவம்}}, \emph{tattuvam}.) An ideology? (\textbf{\ta{சித்தாந்தம்}}, \emph{cittāntam}.) A policy? (\textbf{\ta{கொள்கை}}, \emph{koḷkai}.) To establish? (\textbf{\ta{நிறுவு}}, \emph{niṟuvu}.) To twist? (\textbf{\ta{திரி}}, \emph{tiri}.) To exaggerate? (\textbf{\ta{மிகைப்படுத்து}}, \emph{mikaippaṭuttu}.) To contradict? (\textbf{\ta{முரண்படு}}, \emph{muraṇpaṭu}.) Uncertainty? (\textbf{\ta{ஐயம்}}, \emph{aiyam}.) To reject? (\textbf{\ta{நிராகரி}}, \emph{nirākari}.) Fairness? (\textbf{\ta{நியாயம்}}, \emph{niyāyam}.) Justice? (\textbf{\ta{நீதி}}, \emph{nīti}.) Ethics? (\textbf{\ta{அறம்}}, \emph{aṟam}.) Blame? (\textbf{\ta{பழி}}, \emph{paḻi}.) An excuse? (\textbf{\ta{சாக்கு}}, \emph{cākku}.) To determine? (\textbf{\ta{தீர்மானி}}, \emph{tīrmāṉi}.) To look into? (\textbf{\ta{பரிசீலி}}, \emph{paricīli}.) A yardstick? (\textbf{\ta{அளவுகோல்}}, \emph{aḷavukōl}.) A good match? (\textbf{\ta{பொருத்தம்}}, \emph{poruttam}.) A connection? (\textbf{\ta{சம்பந்தம்}}, \emph{campantam}.) Research? (\textbf{\ta{ஆராய்ச்சி}}, \emph{ārāycci}.) A study? (\textbf{\ta{ஆய்வு}}, \emph{āyvu}.) Data? (\textbf{\ta{தரவு}}, \emph{taravu}.) An analysis? (\textbf{\ta{பகுப்பாய்வு}}, \emph{pakuppāyvu}.) To analyse closely? (\textbf{\ta{அலசு}}, \emph{alacu}.)
\end{tcolorbox}
